\section*{ADR-005: IPC Protocol Selection}
\addcontentsline{toc}{section}{ADR-005: IPC Protocol Selection}

\begin{description}[leftmargin=3cm,labelwidth=2.5cm]
\item[\textbf{Status}] Accepted
\item[\textbf{Date}] May 2024
\item[\textbf{Authors}] Symphony Communication Team
\end{description}

\subsection*{Summary}

We selected a multi-transport IPC system using Unix domain sockets (primary), shared memory (high-throughput), and named pipes (Windows compatibility) for communication between the Conductor and out-of-process extensions, rather than network-based protocols or single-transport solutions.

\subsection*{Context}

Symphony's IPC system must handle:
\begin{itemize}
\item Low-latency communication (target <1ms)
\item High-throughput data transfer for AI models
\item Cross-platform compatibility (Windows, macOS, Linux)
\item Security isolation between extensions
\item Reliable message delivery and error handling
\end{itemize}

\subsection*{Decision}

We implemented a multi-transport IPC Bus with:
\begin{itemize}
\item Unix domain sockets for general communication
\item Shared memory for large data transfers
\item Named pipes for Windows-specific scenarios
\item JSON-RPC protocol for structured messaging
\item Automatic transport selection based on message characteristics
\end{itemize}

\subsection*{Rationale}

This approach optimizes for different communication patterns while maintaining cross-platform compatibility and security isolation.

\subsection*{Success Criteria}

\begin{itemize}
\item IPC latency under 1ms (Achieved: 0.1-0.5ms)
\item Cross-platform compatibility (Achieved)
\item Reliable message delivery (Achieved: 99.99\% success rate)
\end{itemize}